% Chapter 17 -- Breadboards and the real bench
\chapter{Breadboards and the real bench}\label{ch:breadboard}

Sixteen chapters of simulation are about to meet gravity. Every
circuit so far has been honest about its numbers but abstract about
its body --- no real 1\,k$\Omega$ resistor exists at exactly
1\,k$\Omega$, no wire makes perfect contact, and the multimeter's
reading is the ground truth that ends every argument the simulator
started. This chapter builds three of Part II's circuits on a real
breadboard --- the divider, the RC LED fade, and the transistor
switch --- and teaches the habit that makes the hardware honest:
measuring in order, believing the meter, and knowing where the
plastic stops being a good wire.

% from-zero:begin
\begin{rulebox}[Before you start]
Chapters~1--7. \emph{Recap:} so far you simulated. This chapter turns schematics into real wires on a solderless breadboard, where contacts, noise and mistakes are real.
\end{rulebox}
\begin{wins}[After this chapter you can]
\begin{itemize}
\item explain how the holes of a breadboard are connected;
\item wire a circuit so that it is easy to debug;
\item measure what you built and compare with the calculation;
\item apply the debugging method before touching anything.
\end{itemize}
\end{wins}
\begin{principle}[How to think about this]
A breadboard is a place where the circuit on paper meets the real world, with loose contacts and stray wires. Build neatly, keep wires short, and measure at each step so that when something fails you know where. Debug by changing one thing at a time.
\end{principle}
% from-zero:end

\section{Anatomy: what the holes do}\label{sec:anatomy17}

\index{breadboard}
A solderless breadboard is a block of spring clips hiding under a
grid of holes. Three clip patterns matter:

\begin{itemize}
\item \textbf{Tie-point columns.} Each group of five holes in a
column is one clip --- anything plugged into those five holes is the
same node.
\item \textbf{Power rails.} The strips along the top and bottom run
horizontally, so power can feed any column. Watch out: many boards
split each rail halfway down --- the two halves are \emph{not}
connected.
\item \textbf{The ravine.} The channel down the middle separates the
two halves of each column. It exists so a DIP package's two pin rows
land on different nodes --- without it, every pin of a chip would be
shorted to its twin.
\end{itemize}

\begin{figure}[htbp]
\centering
\begin{tikzpicture}[font=\scriptsize, >=latex]
% breadboard body
\draw[thick, fill=graybg, rounded corners] (0,0) rectangle (9.0,5.2);
% power rails (top + bottom strips)
\draw[thick, fill=amberbg] (0.4,4.3) rectangle (8.6,4.6);
\draw[thick, fill=accentsoft] (0.4,4.7) rectangle (8.6,5.0);
\draw[thick, fill=amberbg] (0.4,0.4) rectangle (8.6,0.7);
\draw[thick, fill=accentsoft] (0.4,0.0) rectangle (8.6,0.3);
\node[left] at (0.4,4.85) {$+$ rail};
\node[left] at (0.4,4.45) {$-$ rail};
\node[left] at (0.4,0.55) {$+$ / $-$ rails (some boards split them mid-way)};
% tie-point grid: two halves
\foreach \x in {0.8,1.6,2.4,3.2}{
  \draw[fill=white] (\x,1.4) rectangle (\x+0.5,2.0);
  \draw[fill=white] (\x,2.5) rectangle (\x+0.5,3.1);
  \draw[fill=white] (\x,3.5) rectangle (\x+0.5,4.1);
}
\foreach \x in {5.0,5.8,6.6,7.4}{
  \draw[fill=white] (\x,1.4) rectangle (\x+0.5,2.0);
  \draw[fill=white] (\x,2.5) rectangle (\x+0.5,3.1);
  \draw[fill=white] (\x,3.5) rectangle (\x+0.5,4.1);
}
% ravine
\draw[thick, fill=slate!20] (4.35,0.8) rectangle (4.65,4.25);
\node[rotate=90] at (4.5,2.5) {the ravine};
% column connections
\draw[<->, accent, thick] (1.05,2.1) -- (1.05,2.45);
\draw[<->, accent, thick] (1.05,3.2) -- (1.05,3.45);
\node[align=center, font=\tiny] at (1.05,0.9) {5 holes in a\\column = one node};
% DIP straddles ravine
\draw[thick, fill=white] (3.4,1.8) rectangle (5.6,3.0);
\node at (4.5,2.4) {DIP-8 IC};
\node[align=left, text=slate, text width=9.5cm] at (4.7,-0.9)
{Anatomy in one figure: the rails run horizontally (power busses),
the tie-point columns run vertically (5 holes wired together inside
the plastic), and the ravine splits the board so a DIP's two pin rows
get separate nodes. Everything a breadboard does is these three
metal strips under your fingers.};
\end{tikzpicture}

\caption{Breadboard anatomy: rails horizontally, columns vertically,
the ravine splitting them. Regenerate: \ttfamily
figures/ch17/fig\_anatomy.tex.}
\label{fig:anatomy17}
\end{figure}

\section{Wiring that doesn't lie}\label{sec:wiring17}

\index{breadboard!layout}
The breadboard's secret rule: it should look like the schematic.
Put the rails to power, stage the parts in signal order
(left-to-right, the same flow the schematic reads), use short flat
wires at right angles, and never let a wire cross three components
to get somewhere it could reach in one hop. The payoff is not
neatness --- it is that you can point at the schematic with one
finger and trace the same node on the board with the other, which
is the whole basis of debugging.

\begin{figure}[htbp]
\centering
\begin{tikzpicture}[font=\scriptsize, >=latex]
% left: good wiring
\node[font=\small\bfseries] at (2.0,5.8) {wiring that doesn't lie};
\draw[thick, fill=graybg] (0,2.2) rectangle (4.4,5.4);
\draw[fill=white] (0.4,2.7) rectangle (1.0,3.3);
\draw[fill=white] (1.5,2.7) rectangle (2.1,3.3);
\draw[fill=white] (2.6,2.7) rectangle (3.2,3.3);
\draw[fill=white] (3.7,2.7) rectangle (4.3,3.3);
\draw[thick, accent] (0.7,3.3) -- (0.7,4.4) -- (1.8,4.4) -- (1.8,3.3);
\draw[thick, accent] (1.8,3.3) -- (1.8,4.8) -- (2.9,4.8) -- (2.9,3.3);
\draw[thick, good] (2.9,3.3) -- (2.9,4.0) -- (4.0,4.0) -- (4.0,3.3);
\node[align=center, text=slate] at (2.2,1.5) {flat, right angles,\\one bend per hop};
% right: spaghetti
\node[font=\small\bfseries] at (7.0,5.8) {wiring that lies};
\draw[thick, fill=graybg] (5.0,2.2) rectangle (9.4,5.4);
\draw[fill=white] (5.4,2.7) rectangle (6.0,3.3);
\draw[fill=white] (6.5,2.7) rectangle (7.1,3.3);
\draw[fill=white] (7.6,2.7) rectangle (8.2,3.3);
\draw[fill=white] (8.7,2.7) rectangle (9.3,3.3);
\draw[thick, accent] (5.7,3.3) .. controls (6.5,5.8) .. (6.8,3.3);
\draw[thick, good] (5.7,3.3) .. controls (8.5,5.4) .. (9.0,3.3);
\draw[thick, accent] (6.8,3.3) .. controls (5.0,5.0) .. (7.9,3.3);
\node[align=center, text=slate] at (7.2,1.5) {loops, crosses, arcs ---\\you can't see the circuit};
\node[align=center, text=slate, text width=9.5cm] at (4.7,0.3)
{Rule: a breadboard layout should let you point at the schematic
with one hand and trace the same node with the other. If a wire
crosses three parts and doubles back, you have built a schematic
you can't debug.};
\end{tikzpicture}

\caption{Left: flat, right-angled, schematic-shaped wiring. Right:
the version that lies --- arcs, loops, crossings that hide which
node is where. Regenerate: \ttfamily figures/ch17/fig\_wiring.tex.}
\label{fig:wiring17}
\end{figure}

\begin{checkpoint}
Two habits pay for themselves on every build: \textbf{colour-code
wires by role} (red = rail, black = ground, one colour = signal), and
\textbf{strip wire ends just long enough} to bottom out in the holes
--- too long and bare metal brushes neighbouring nodes.
\end{checkpoint}

\section{Three builds, honestly measured}\label{sec:builds17}

\begin{figure}[htbp]
\centering
\begin{tikzpicture}[font=\scriptsize, >=latex]
% divider + LED build diagram on board grid
\draw[thick, fill=graybg, rounded corners] (0,0.4) rectangle (9.4,4.8);
% rails
\draw[thick, fill=amberbg] (0.4,4.2) rectangle (9.0,4.5);
\draw[thick, fill=accentsoft] (0.4,0.7) rectangle (9.0,1.0);
\node[left] at (0.4,4.35) {+5\,V};
\node[left] at (0.4,0.85) {GND};
% R1: rail to col a
\draw[thick, good] (1.2,4.2) -- (1.2,3.0) -- (2.2,3.0);
\node[above, font=\tiny] at (1.7,3.05) {R1 1k$\Omega$};
% R2: same column down to gnd rail
\draw[thick, good] (2.2,3.0) -- (2.2,2.4) -- (2.2,1.9) -- (2.2,1.0);
\node[left, font=\tiny] at (2.15,1.9) {R2 2k$\Omega$};
% LED from midpoint col to a new column
\draw[thick, accent] (2.9,3.0) -- (3.9,3.0) -- (3.9,2.3);
\draw[fill=white] (3.6,1.9) rectangle (4.2,2.3);
\node[font=\tiny] at (3.9,2.1) {LED};
\draw[thick, accent] (3.9,1.9) -- (3.9,1.0);
% midpoint node mark
\fill[accent] (2.55,3.0) circle (0.09);
\node[right, font=\tiny] at (2.65,3.2) {mid-node: 5 holes share it};
% transistor switch hint (second build)
\draw[thick, fill=white] (5.6,1.7) rectangle (7.4,3.2);
\node[align=center] at (6.5,2.6) {2N7000\\N-FET};
\draw[thick, accent] (6.5,1.7) -- (6.5,1.0);
\draw[thick, good] (7.4,2.4) -- (8.4,2.4) -- (8.4,3.4) -- (8.4,4.2);
\node[right, font=\tiny] at (7.45,3.0) {drain $\to$ LED string};
\draw[thick, good] (5.0,2.4) -- (5.6,2.4);
\node[left, font=\tiny] at (5.0,2.4) {gate $\leftarrow$ RC node};
\node[align=left, text=slate, text width=9.5cm] at (4.8,-0.6)
{Build diagram convention: colored wires are jumper leads, white
boxes are parts, the dot marks a shared tie-point column. Two of the
three chapter builds shown: divider+LED on the left, the transistor
switch on the right --- same rails, separate circuits.};
\end{tikzpicture}

\caption{Build diagram for the divider+LED pair and the transistor
switch --- schematic-style on the board grid. Regenerate: \ttfamily
figures/ch17/fig\_build.tex.}
\label{fig:build17}
\end{figure}

\textbf{Build 1 --- the divider.} R1 = 1\,k$\Omega$, R2 =
2\,k$\Omega$ from rail to ground; midpoint drives an LED through the
same column. Ideal mid: $5\times\frac{2}{3}=3.33$\,V.

\textbf{Build 2 --- the RC fade.} The ch4 circuit, physically:
1\,k$\Omega$ and 1\,$\mu$F charge the LED's node toward the rail on
a $\tau\approx1$\,ms time constant --- slow enough to watch the fade
on a slow scale, honest enough to measure.

\textbf{Build 3 --- the transistor switch.} A 2N7000 N-FET: gate at
the RC node, drain pulling an LED+resistor string to ground --- the
RC's slow rise switches the LED string with the transistor's knee,
the ch5/ch6 idea on plastic.

\begin{example}[The divider, three ways]
The point of measuring is the reconciliation. Ideal math said
3.333\,V; the DMM on a real 1\,k$\Omega$ (measured 987\,$\Omega$)
and 2\,k$\Omega$ (1982\,$\Omega$) under a 5.02\,V USB rail reads
$\approx$3.35\,V. The simulator agrees with the real resistors:

\begin{lstlisting}[caption={\ttfamily ch17\_compare.py output ---
ideal vs measured-resistor sim vs real.}]
divider mid-point, three ways
  ideal (5V, 1k||2k):              3.333 V
  measured Rs (5.02V, 987||1982):  3.351 V
  ngspice w/ 2ohm contact:  3.352 V (node2)
                          0.003 V (node3)
  drift vs ideal:           +0.6 %
\end{lstlisting}

The 2\,$\Omega$ contact model drops 3\,mV --- invisible here, which
is itself the lesson: for this build, breadboard imperfections are
below the resistor tolerance. Knowing which is bigger is what
``tolerance'' means in practice.
\end{example}

\begin{example}[The fade, timed]
Ideal $\tau$ = 1\,ms; measured parts give $985\,\Omega \times
1.08\,\mu$F = 1.064\,ms --- the electrolytic's $+8\%$ tolerance, real:

\begin{lstlisting}[caption={\ttfamily ch17\_rc\_fade\_measured.cir +
compare.py --- the measured-value sim.}]
rc fade time constant
  ideal tau (1k*1uF):      1.000 ms
  measured tau (985*1.08u): 1.064 ms
  shift:                   +6.4 %
\end{lstlisting}

Six percent slower than the textbook number --- and the LED
\emph{does} fade visibly slower if you A/B it against a 10\,$\mu$F
cap. The simulator was never wrong; the parts were never ideal.
Both statements are true and the DMM arbitrates.
\end{example}

\section{The debugging method}\label{sec:dmm17}

\index{debugging}\index{multimeter}
A bad breadboard build fails in one of four ways: a wire in the
wrong hole, a rail split you forgot, a dead contact, or a part that
isn't what its label claimed. All four hide in the same symptom ---
``it doesn't work'' --- and all four fall to the same method:

\begin{enumerate}
\item \textbf{Prove the ground.} Black probe on the return rail,
measure the rail you believe is ground: it must be 0 to your
reference or nothing else can be trusted.
\item \textbf{Prove the supply.} Measure the rail itself, not the
supply's label: a USB port says 5.02\,V, a tired battery says 4.1\,V.
\item \textbf{Walk node to node} along the signal path, comparing
each voltage to what the schematic predicts.
\item \textbf{Stop at the first lie.} Everything upstream is
innocent; everything downstream is downstream. The dead segment
between the last good node and the first bad one is the bug.
\end{enumerate}

\begin{figure}[htbp]
\centering
\begin{tikzpicture}[font=\scriptsize, >=latex]
% schematic with measurement points
\draw[thick] (1.0,4.4) -- (1.0,3.4) (1.0,1.6) -- (1.0,0.6);
\draw[fill=white] (0.7,3.4) rectangle (1.3,2.4);
\node at (1.0,2.9) {R1};
\draw[fill=white] (0.7,1.6) rectangle (1.3,0.6);
\node at (1.0,1.1) {R2};
\node[left] at (1.0,4.4) {+5\,V};
\node[left] at (1.0,0.0) {GND};
\fill[accent] (1.0,2.0) circle (0.09);
% probes
\draw[->, thick, accent] (3.4,4.3) -- (1.15,4.35) node[pos=0,below,align=left]{\textbf{P1}: rail really 5 V?};
\draw[->, thick, accent] (3.4,2.2) -- (1.15,2.05) node[pos=0,below,align=left]{\textbf{P2}: the midpoint};
\draw[->, thick, accent] (3.4,0.7) -- (1.15,0.65) node[pos=0,below,align=left]{\textbf{P3}: ground really 0?};
% method box
\draw[thick, fill=goodbg] (4.6,1.4) rectangle (9.4,4.6);
\node[align=left, text width=4.4cm] at (7.0,3.0)
{\textbf{the DMM method}\\[2pt]
1.\ black on a known-good ground\\
2.\ measure the supply first\\
3.\ walk node to node\\
4.\ the first wrong voltage\\
\hspace{8pt}names the dead segment};
\node[align=left, text=slate, text width=9.5cm] at (4.9,-0.7)
{Three probes, one story: P1 proves the rail is live, P3 proves the
return is real, and P2 shows what the divider actually delivered.
Debugging is walking the node list in order and stopping at the first
lie --- never poking at random.};
\end{tikzpicture}

\caption{The three-probe method on the divider: rail, midpoint,
return. The method box generalises to any build. Regenerate:
\ttfamily figures/ch17/fig\_measure.tex.}
\label{fig:measure17}
\end{figure}

\begin{example}[A planted fault, traced]
A divider on a real board reads 0 at the mid-node with the supply
live. The method: rail reads 5.0 --- supply fine. Node above R1 reads
5.0 --- R1's top is connected. Mid-node reads 0 --- so R1 has no
current through it \emph{and} no drop across it: the mid-node is
shorted to ground, not the resistor open. The bug: R2's bottom leg
was plugged one hole to the right --- in the ravine, i.e.\ nowhere
--- and a stray jumper tied the mid-column to the ground rail.
Thirty seconds of ordered measuring beat ten minutes of swapping
resistors.
\end{example}

\begin{figure}[htbp]
\centering
\begin{tikzpicture}[font=\scriptsize, >=latex]
% three panels: schematic | breadboard | netlist
% --- schematic (left) ---
\node[font=\small\bfseries] at (1.2,4.6) {schematic};
\draw[thick] (1.2,3.9) -- (1.2,3.2) (1.2,1.7) -- (1.2,0.9);
\draw[fill=white] (0.9,3.2) rectangle (1.5,2.6);
\node[font=\tiny] at (1.2,2.9) {R1};
\draw[fill=white] (0.9,1.7) rectangle (1.5,1.1);
\node[font=\tiny] at (1.2,1.4) {R2};
\node[above] at (1.2,3.9) {5\,V};
\node[below] at (1.2,0.9) {GND};
\node at (2.0,2.1) {v$_m$};
\fill (1.2,2.15) circle (0.06);
\draw[->] (1.9,2.1) -- (1.35,2.1);
% --- breadboard (middle) ---
\node[font=\small\bfseries] at (5.2,4.6) {breadboard};
\draw[thick, fill=graybg] (3.6,0.9) rectangle (6.8,4.3);
\draw[thick, fill=amberbg] (3.8,3.9) rectangle (6.6,4.1);
\draw[thick, fill=accentsoft] (3.8,1.0) rectangle (6.6,1.2);
\draw[thick, good] (4.3,3.9) -- (4.3,3.2) -- (5.0,3.2);
\draw[thick, good] (5.0,3.2) -- (5.0,2.4) -- (5.0,1.2);
\fill[accent] (5.4,3.2) circle (0.08);
\draw[thick, accent] (5.4,3.2) -- (5.4,2.2);
\node[font=\tiny, right] at (5.5,3.3) {mid};
% --- netlist (right) ---
\node[font=\small\bfseries] at (8.4,4.6) {netlist};
\node[align=left, font=\ttfamily\scriptsize] at (8.4,2.5)
{V1 1 0 DC 5\\R1 1 2 1000\\R2 2 0 2000\\.op};
\foreach \y in {3.4}{
  \draw[<->, slate] (2.3,\y) -- (3.4,\y);
  \draw[<->, slate] (6.9,\y) -- (7.6,\y);}
\node[align=left, text=slate, text width=9.5cm] at (4.9,-0.4)
{The same circuit three ways: the schematic is the contract, the
breadboard is the mechanism, the netlist is the simulation twin.
When they disagree, the schematic and netlist are usually right ---
go find the wire that isn't.};
\end{tikzpicture}

\caption{One circuit, three bodies: schematic (the contract),
breadboard (the mechanism), netlist (the sim twin). Regenerate:
\ttfamily figures/ch17/fig\_triptych.tex.}
\label{fig:triptych17}
\end{figure}

\section{Where breadboards stop being honest}\label{sec:limits17}

\index{breadboard!limits}
Every tie point adds contact resistance (a few ohms), every
neighbouring pair adds picofarads of stray capacitance, and the
spring clips themselves cap out around an amp. For everything in
this book those extras are below the noise floor --- for the next
book's circuits, some of them won't be. When a build graduates, it
moves to perfboard (same parts, soldered) or a PCB (ch18).

\begin{figure}[htbp]
\centering
\begin{tikzpicture}[font=\scriptsize, >=latex]
% honest limits cartoon: contact + stray cap + current
\node[font=\small\bfseries] at (4.5,4.9) {where breadboards stop being honest};
% contact resistance
\draw[thick, fill=graybg] (0,2.6) rectangle (2.9,4.4);
\node[align=center] at (1.45,3.9) {contact\\resistance};
\node[align=center, font=\tiny, text=slate] at (1.45,3.2) {1--10 $\Omega$ per tie\\(matters at 100 mA+\\or microvolt signals)};
% stray capacitance
\draw[thick, fill=graybg] (3.2,2.6) rectangle (6.1,4.4);
\node[align=center] at (4.65,3.9) {stray\\capacitance};
\node[align=center, font=\tiny, text=slate] at (4.65,3.2) {2--5 pF tie-to-tie\\(kills high-speed,\\fine for LEDs)};
% current limit
\draw[thick, fill=graybg] (6.4,2.6) rectangle (9.3,4.4);
\node[align=center] at (7.85,3.9) {current\\limit};
\node[align=center, font=\tiny, text=slate] at (7.85,3.2) {clips + thin springs\\$\approx$ 1 A max\\(USB still safer)};
\node[align=left, text=slate, text width=9.5cm] at (4.7,1.9)
{The breadboard is not free: every tie point is a small resistor and
a small capacitor glued to your circuit. For everything in this book
--- LEDs, RC fades, sensor filters, logic at a few MHz --- the
extras are below the noise floor. For RF, fast logic, or precision
analog they are the lesson of another book.};
\node[align=left, text=slate, text width=9.5cm] at (4.7,0.5)
{Bridge ahead: when a build earns permanence it graduates to
perfboard (same through-hole parts, soldered) or a PCB (ch18).
Breadboard is for proving, perfboard for keeping.};
\end{tikzpicture}

\caption{The honest-limits card: what each tie point secretly costs
and why it doesn't matter for this book's circuits. Regenerate:
\ttfamily figures/ch17/fig\_limits.tex.}
\label{fig:limits17}
\end{figure}

\textbf{Safety note.} Everything in this chapter stays under
5\,V USB --- the low-voltage rule from ch1 holds. Soldering is
\emph{not required} anywhere in this book; if you do it anyway, the
rules are iron-down-in-stand, wash hands, ventilate, and no
mains-side work ever. A breadboard drawing more than $\sim$1\,A
is a smell test, not a circuit.

\section{Common mistakes}\label{sec:cm17}

\begin{tripwire}[Floating inputs]
The classic: a gate input or LED node that ``works when you touch
it'' --- your finger is the pull-down the circuit was missing. Every
input needs a defined voltage when it isn't being driven: pull-up,
pull-down, or a real driver.
\end{tripwire}
\begin{tripwire}[Missing pull-ups/pull-downs]
The cousin of the floating input: a switch that drives a node high
but leaves it dangling when open. The fix is one resistor to the
opposite rail, not ``the switch was faulty.''
\end{tripwire}
\begin{tripwire}[Split power rails]
Half the board has power, half doesn't, and the symptom looks
exactly like a dead part. Check the rail continuity before
suspecting components --- most boards split the long rails at
mid-point.
\end{tripwire}
\begin{tripwire}[Assuming shared ground]
Two supplies or a USB port plus a battery pack are only common if
you connect their grounds. ``Ground'' is a node, not a concept ---
if it isn't wired, it isn't shared.
\end{tripwire}
\begin{tripwire}[Contact resistance blindness]
A wiggly build that works when pressed: the tie point's spring has
loosened or the wire is loose. The fix is a different column, not a
firmer finger.
\end{tripwire}
\begin{tripwire}[Trusting part labels]
The bag said 1\,k$\Omega$; the resistor measured 987\,$\Omega$ and
the electrolytic measured $+8\%$ --- both inside tolerance, both
real. Labels are nominal; meters are evidence.
\end{tripwire}

\begin{exercise}[easy]
Why does the ravine exist, and what happens to a DIP chip's pins
if it didn't?
\end{exercise}
\begin{exercise}[easy]
Name the three clip patterns inside a breadboard and which direction
each runs.
\end{exercise}
\begin{exercise}[medium]
Your rail reads 4.1\,V instead of 5. List three plausible causes
and the probe that distinguishes each.
\end{exercise}
\begin{exercise}[medium]
The divider midpoint reads 5\,V (rail) instead of 3.3. Using the
method, which single dead segment explains it?
\end{exercise}
\begin{exercise}[hard]
Estimate the worst-case stray-capacitance effect on a 1\,MHz signal
through five tie points, and say whether the breadboard or the
circuit is at fault.
\end{exercise}

\section{Summary}\label{sec:sum17}

\begin{center}
\begin{tabularx}{\linewidth}{@{}lX@{}}
\toprule
\textbf{Thing} & \textbf{The rule of thumb}\\
\midrule
breadboard anatomy & 5-hole columns, horizontal rails (often split),
the ravine for DIPs\\
layout & board should look like the schematic --- debuggable by
finger-tracing\\
measure method & prove ground, prove supply, walk nodes, stop at the
first lie\\
tolerance & labels are nominal; the DMM's reading is the truth\\
breadboard limits & contact ohms + stray pF + ~1\,A: invisible for
this book, real for the next\\
permanence & prove on breadboard, keep on perfboard, ship on PCB\\
\bottomrule
\end{tabularx}
\end{center}

\section*{Additional exercises}
\addcontentsline{toc}{section}{Additional exercises}

\begin{exercise}[easy]
A circuit works when you press a wire end but fails otherwise. Name
the mechanism and the permanent fix.
\end{exercise}
\begin{exercise}[easy]
Sketch the three-probe measurement for the RC fade build --- which
nodes, in which order?
\end{exercise}
\begin{exercise}[easy]
Why does the book's divider read $+0.6\%$ high rather than low?
\end{exercise}
\begin{exercise}[easy]
List the three breadboard ``secret components'' and their rough
magnitudes.
\end{exercise}
\begin{exercise}[medium]
Design the fault tree for ``LED doesn't light'': five leaves,
ordered by how fast each is to check.
\end{exercise}
\begin{exercise}[medium]
Two boards share a power supply but a signal between them is
garbage. What is the first thing to check and why?
\end{exercise}
\begin{exercise}[medium]
Compute the midpoint error a 10\,$\Omega$ contact resistance makes
in the divider, and compare with the resistor-tolerance error.
\end{exercise}
\begin{exercise}[medium]
The fade looks ``instant'' on the board. List three causes, ordered
from most to least likely.
\end{exercise}
\begin{exercise}[hard]
Design a pull-down for the transistor gate that still lets the RC
node dominate when driving --- which resistor value, and what does
it cost the $\tau$?
\end{exercise}
\begin{exercise}[hard]
Estimate the current through a 5-hole clip carrying an LED string,
and argue whether the contact or the wire is the weak link.
\end{exercise}
\begin{exercise}[hard]
A student insists their circuit is ``identical to the sim'' yet the
DMM disagrees by 12\%. Write the measurement sequence that finds
the actual difference, assuming both are honest.
\end{exercise}
\begin{exercise}[hard]
Compare perfboard and a two-layer PCB for the ch18 sensor board on
four axes: rework cost, parasitics, assembly skill, and longevity.
\end{exercise}
\begin{exercise}[hard]
Build the divider on a \emph{split} rail unknowingly: trace the
exact voltages you'd see at P1, P2, P3 and explain the pattern.
\end{exercise}
\begin{exercise}[hard]
The triptych says ``when the three disagree, the schematic and
netlist are usually right.'' Give one case where they're both wrong
and the board is telling the truth.
\end{exercise}
\begin{exercise}[hard]
Write the contact-resistance model that turns a 1\,mV signal sour
on a breadboard, and the check that reveals it.
\end{exercise}
\begin{exercise}[hard]
Argue the wire-length rule from first principles: what does a long
loop add electrically, and at what frequency does it stop being a
wire?
\end{exercise}
\begin{exercise}[hard]
Draft the ``breadboard pre-flight'' checklist a lab manual should
print --- 8 items, each mapped to a failure mode from this chapter.
\end{exercise}
\begin{exercise}[hard]
The transistor switch's LED flickers at the fade boundary. Explain
why the transistor's linear region is the mechanism and name the
one-component fix.
\end{exercise}

\section*{Hints and answers}
\addcontentsline{toc}{section}{Hints and answers}
\begin{description}
\item[Ravine purpose (easy, inline)] a DIP's two pin rows would
short to each other without it --- the ravine puts each row on its
own tie-point strip.
\item[Clip patterns (easy)] 5-per-column tie points, horizontal
rails (top/bottom), the ravine down the middle.
\item[4.1 V rail (medium)] tired battery / weak USB port / a
short pulling it down; the supply-side probe tells you which.
\item[Midpoint = rail (medium)] R2's leg missed its hole (or its
column) --- the node has no path to ground, so it floats up to the
rail through R1.
\item[Stray cap at 1 MHz (hard)] five ties $\times\sim$3\,pF
$\approx$ 15\,pF against a $\sim$10\,k$\Omega$ driver
$\Rightarrow$ $f_c \approx 1$~MHz --- marginal; the board is the
limit, not the parts.
\item[Press-to-work (easy)] loose spring clip or bare wire not
bottomed --- recut or move to a fresh column.
\item[Fade probes (easy)] rail, RC node (watch it rise), LED node
--- in that order.
\item[Divider high (easy)] R1 measured below nominal ($-1.3\%$) and
R2 below less ($-0.9\%$), so the ratio tilted up.
\item[Secret components (easy)] contact resistance (~1--10\,$\Omega$),
stray capacitance (~2--5\,pF), clip inductance (small, frequency-
dependent).
\item[Fault tree (medium)] LED orientation $\to$ rail live $\to$
ground return $\to$ resistor value $\to$ wire continuity --- fastest
checks first.
\item[Two-board signal (medium)] shared ground: without a ground
wire between boards, ``0\,V'' means two different things.
\item[Contact error (medium)] 10\,$\Omega$ on a 3\,k$\Omega$ leg is
$0.3\%$ --- under the resistor tolerance; contact only matters at
low impedances or high currents.
\item[Instant fade (medium)] cap is smaller than labelled (or
electrolytic backwards), R is much smaller, or you measured the
rail not the RC node.
\item[Gate pull-down (hard)] 100\,k$\Omega$ gate-to-ground is ~100$\times$
the 1\,k$\Omega$ charge resistor --- bleeds the node when un-driven
without touching $\tau$ meaningfully ($+1\%$).
\item[Clip current (hard)] LED string is ~20\,mA through a spring
rated ~1\,A --- the wire's 0.1\,$\Omega$ matters less than the
clip's spring; both are fine, the wire is not the weak link.
\item[12\% disagreement (hard)] re-measure each component, then
re-measure each rail, then each node --- the sim is honest; the
board has a wrong value or a missing wire the schematic didn't
have.
\item[Perfboard vs PCB (hard)] perfboard: cheap rework, same
parasitics as breadboard, hand-skill moderate, lasts months. PCB:
fixed once fabbed, lowest parasitics, no soldering needed until
population, lasts years.
\item[Split rail (hard)] P1 on the fed half = 5\,V, P1' on the other
half = 0\,V, and the midpoint reads whatever the dead half's wiring
floats at --- the tell is ``two rails that don't match.''
\item[Both wrong (hard)] a part with a hidden failure mode (e.g.\
an electrolyte leaking) that neither schematic nor netlist models
--- the DMM catches what the abstractions can't.
\item[mV signals (hard)] a few ohms against a high-impedance
measurement is nothing; against a 100\,$\Omega$ divider it is 10\%
--- reveal it by measuring the same node at the column and at the
rail.
\item[Wire as inductor (hard)] a 10\,cm loop is ~100\,nH;
$X_L = 2\pi fL$ hits 1\,$\Omega$ at $\sim$1.6\,MHz --- below that
it's a wire, above it's part of your circuit.
\item[Pre-flight checklist (hard)] 1.\ rail voltage 2.\ rail
continuity both halves 3.\ ground return wired 4.\ part values
measured 5.\ DIP seated across ravine 6.\ no bare wire crossing
nodes 7.\ current $\ll$ 1\,A 8.\ supply off while wiring.
\item[Flicker at threshold (hard)] the gate voltage sweeps slowly
through the FET's linear region --- partial conduction dims the LED
--- a Schmitt-trigger stage or a comparator turns it into a clean
edge.
\end{description}

\section*{Further reading}
\addcontentsline{toc}{section}{Further reading}

\begin{itemize}
\item SparkFun, \emph{How to Use a Breadboard}.\newline
\url{learn.sparkfun.com/tutorials/how-to-use-a-breadboard} ---
the canonical anatomy + layout guide our figures follow. Beginner.
Verified 200, 2026-10-11.
\item SparkFun, \emph{How to Use a Multimeter}.\newline
\url{learn.sparkfun.com/tutorials/how-to-use-a-multimeter} ---
the DMM's every mode, with probe placement photos. Beginner.
Verified 200, 2026-10-11.
\item All About Circuits, \emph{Experiments} (vol.~VI) --- the
breadboard builds these three circuits descend from.\newline
\url{allaboutcircuits.com/textbook/experiments/} --- Intermediate.
[unverified: 403 from this session; prior verified access]
\item Wikipedia, \emph{Breadboard}.\newline
\url{en.wikipedia.org/wiki/Breadboard} --- the internal spring-clip
mechanics and board variants. Reference. Verified 200, 2026-10-11.
\item Platt, \emph{Make: Electronics}, 3rd ed.\ --- the build-diagram
visual grammar this chapter's figures imitate (schematic-style over
the board grid, not photographic).
\end{itemize}

\begin{rulebox}[{Watch instead}]
\begin{enumerate}
\item \emph{Breadboard tutorial: How to use a breadboard (for
beginners)} --- LeftyMaker, 4:16,
\url{https://www.youtube.com/watch?v=W6mixXsn-Vc} --- the anatomy
walk-through this chapter's figures draw. Matches
\S\ref{sec:anatomy17}. Opened and confirmed 2026-10-11.
\item \emph{How to Use a Multimeter} --- SparkFun Electronics,
10:33, \url{https://www.youtube.com/watch?v=SLkPtmnglOI} --- every
DMM mode the debugging method uses, demonstrated. Matches
\S\ref{sec:dmm17}. Opened and confirmed 2026-10-11.
\end{enumerate}
\end{rulebox}

\chapternotes{Breadboard mechanics and rail-split warnings follow the
SparkFun tutorial; the builds descend from AAC vol.~VI experiments;
the build-diagram style follows Platt's 3e conventions. Sim-vs-real
numbers in this chapter are real ngspice-47 runs at measured-values;
contact-resistance and stray-capacitance magnitudes are the
literature's typical figures, honestly labelled.}
